\section{The full-dimensional seed and the two propagation runs}
\label{sec:actual-seed-two-propagations}

We continue with the actual arithmetic frame, entry spaces,
and ordinary Hecke ideals constructed above.
For \(n=2,3\) write
\[
 \mathcal B_n\twoheadrightarrow
       \mathbb T_n=\mathcal B_n/J_n,\qquad
 \Phi_n:\mathcal B_n\twoheadrightarrow B_n,\qquad
 M_n=J_nB_n .
                                                               \tag{SF1}
\]
The rings \(A_n,B_n\), their grading, and \(\Phi_n\)
are the rings and maps of
Sections~\ref{sec:entry-space-flat} and
\ref{sec:det-entry-finite-corner}.
The modules defining \(\mathbb T_n\), including their
tame coefficient lattices, are the modules of
Section~\ref{sec:common-ordinary-arithmetic-action}.
We use the two actual cuspidal seeds, their symmetric
conjugate-self-dual pairings, and their insertion in
these same ordinary modules.
This section turns those points into full-dimensional
seeds and then performs the two propagation runs in
their required order.

We fix the exact weight augmentation notation.
Put \(k=\mathcal O/\mathfrak m_{\mathcal O}\) and \(q=|k|\).
The ideal used by \srcid{Space.FinLam} is
\[
 \mathfrak l_n=\mathfrak m_{\mathcal O}B_n+
   \bigl(\psi_{vi}(g)-1:
       g\in I_v,\ \Psi(\operatorname{loc}_v g)=1\bigr).
                                                               \tag{SF2}
\]
Curve separation instead uses
\[
 N_v=I_v\cap\ker(\varepsilon\circ\operatorname{loc}_v).
\]
The frame supplies \(N_v\subseteq\ker(\Psi\circ\operatorname{loc}_v)\).
Its \(d_v\) local inertia test elements lie in the latter
kernel; all but one lie in \(N_v\).
These are distinct assertions.

\subsection{Finite weight support and all component dimensions}
Each \(\mathbb T_n\) is finite and torsion-free over its
weight power-series ring
\(\Lambda_n\simeq\mathcal O[[X_1,\ldots,X_{nd}]]\).
It is nonzero because its actual cuspidal seed occurs
in the ordinary module.
Every minimal prime \(Z\) over \(J_n\) satisfies
\[
                      \dim\mathcal B_n/Z=1+nd.              \tag{SF3}
\]
Indeed a minimal prime of a Noetherian ring is an
associated prime. No nonzero element of the domain
\(\Lambda_n\) is a zero divisor on \(\mathbb T_n\),
so the contraction of every such prime to
\(\Lambda_n\) is zero. The corresponding quotient is
finite integral over \(\Lambda_n\), and has its dimension.

Also
\[
                    \mathfrak m_{B_n}
                           \subseteq\sqrt{M_n+\mathfrak l_n}.
                                                               \tag{SF4}
\]
Here is the needed finite-weight argument.
Take a prime \(P\supseteq M_n+\mathfrak l_n\) and pull
it back to \(\mathcal B_n\).
For an inertia element \(g\), the element \(g^{q-1}\)
has trivial \(\Psi\)-value.
Equation (SF2), together with the fixed scalar twist,
therefore makes the \((q-1)\)-st power of each
universal ordinary inertia character equal to one
in the pulled-back quotient.
The scalar factor also has \((q-1)\)-st power one
after reduction modulo \(\mathfrak m_{\mathcal O}\).

A unit \(a\) whose residue is fixed and for which
\(a^{q-1}=1\) is the corresponding constant root of
unity in this local quotient. For two such values
\(a,c\), factor
\[
 a^{q-1}-c^{q-1}
   =(a-c)\sum_{i=0}^{q-2}a^ic^{q-2-i}.
\]
The second factor has nonzero residue
\((q-1)\bar a^{q-2}\), so it is a unit.
The universal property of the inertia weight ring
then makes its maximal ideal vanish in the quotient.
Finiteness over \(\Lambda_n\) makes the pulled-back
prime maximal. Surjectivity of \(\Phi_n\) makes
\(P=\mathfrak m_{B_n}\).
This proves (SF4) by intersecting the primes in question.
It reconstructs \srcid{CurveDVRSpecialization.FinLamOf};
it does not assume the finite-weight conclusion as an
extra property of the entry space.

\subsection{From one actual seed to the generic representation of a component}
After a finite coefficient extension, the actual seed
gives a map
\[
 x_0:\mathcal B_n\longrightarrow\overline{\mathbb Q}_\ell,
 \qquad J_n\subseteq\ker x_0,
                                                               \tag{SF5}
\]
and an absolutely irreducible representation with the
prescribed ordered local characters.
Untwisting by the fixed scalar \(\kappa_n\) gives
the symmetric polarized frame representation constructed
in the seed sections.
It may be normalized as below: its trace and local
characters are integral over one finite coefficient
DVR, its marked difference has the separated residual
eigenvalues, and the simultaneous normalization of
the form and marked difference preserves its ordered
ordinary identities.
Thus (SF5) is the actual normalized seed required
here, rather than only a congruent residual determinant.

Choose a minimal prime \(Z\) over \(J_n\) with
\(Z\subseteq\ker x_0\), and set
\[
 C=\mathcal B_n/Z,\quad
 K=\overline{\operatorname{Frac}(C)},\quad
 p=(D_C,\psi_C).
\]
Then \(C\) is a complete local domain with residue \(k\)
and dimension \(1+nd\), and \(C\hookrightarrow K\)
is injective.
The determinant over \(K\) has a semisimple
representation \(r_0\).
The field-valued determinant theorem used here is
\href{https://gaetan.chenevier.perso.math.cnrs.fr/articles/determinants.pdf}
{Chenevier, Theorem 2.12}; it applies to a determinant
over an algebraically closed field without any
residual irreducibility assumption.

In fact \(r_0\) is absolutely irreducible.
Choose \(n^2\) group matrices spanning the full matrix
algebra at the actual seed. Their trace Gram determinant
is nonzero there. It is the image of
\[
          \Delta=\det\bigl(D_C^{\rm tr}(g_ag_b)\bigr)_{a,b}
                                                        \in C .
                                                               \tag{SF6}
\]
Thus \(\Delta\ne0\) in \(C\), and it remains nonzero
in \(K\). The same Gram matrix therefore makes the
\(n^2\) generic matrices linearly independent.
They span \(\operatorname{Mat}_n(K)\), which proves
irreducibility. This step starts from the actual
irreducible seed; a merely reducible characteristic
polynomial point would not give (SF6).

We give the descent to the frame explicitly, since the
generic matrices over \(K\) have not been equipped with
a topology. Let
\[
 \widetilde r(g)=\kappa_n(g)^{-1}r_0(g),\qquad
 \widetilde t(g)=\kappa_n(g)^{-1}D_C^{\rm tr}(g)\in C .
\]
The scalar twist is integral and continuous.
Full matrix span gives the trace criterion
\[
 \widetilde r(g)=1
 \quad\Longleftrightarrow\quad
 \widetilde t(hg)=\widetilde t(h)\quad\hbox{for every }h .
                                                               \tag{SF7}
\]
For the reverse implication, the trace pairing with
the span of \(\widetilde r(h)\) kills
\(\widetilde r(g)-1\); the matrix trace pairing is
nondegenerate. Since \(\widetilde t\) is \(C\)-adically
continuous and \(C\) is separated, (SF7) makes the
kernel a closed subgroup of the profinite arithmetic
group. It is normal, and the determinant polarization
makes it stable under conjugation by \(c\).

The unramified assertion (CA14) makes this kernel
contain the distinguished inertia subgroup \(I_0\),
where \(\kappa_n=1\).
For an original tame inertia element \(g\) with
\(\Psi(\pi g)=1\), the exact \srcid{TwLink.tame}
condition gives a prime-to-\(\ell\) integer \(m\)
such that
\(\pi(g)^m=\kappa_n(g)^m=1\).
Its \(m\)-th power lies in the residual inertia
kernel killed by the universal pair.
Full matrix span converts that determinant relation
to \(r_0(g^m)=1\).
Thus \(\widetilde r(g)^m=1\), while its integral
characteristic polynomial reduces to \((X-1)^n\).
It follows that \(\widetilde r(g)=1\).
One direct proof writes
\((M-1)(1+M+\cdots+M^{m-1})=0\);
the determinant of the second factor is an integral
polynomial in the characteristic coefficients and
has residue \(m^n\ne0\), so the factor is invertible.

Apply the closed-normal kernel property
\srcid{TwLink.ker}.
It proves that \(\ker\pi\subseteq\ker\widetilde r\).
Surjectivity of \(\pi\) now gives an actual representation
\(r:G_0\to\operatorname{GL}_n(K)\) of the fixed frame group.
The trace descends to a continuous function
\(t:G_0\to C\): the compact-to-Hausdorff surjection
\(\pi\) is a quotient map, and (SF7) makes the descended
trace well defined. The local group isomorphisms in
\srcid{TwLink} similarly give continuous characters
\(\chi_{vi}:G_v\to C^\times\).
They satisfy
\[
 \begin{aligned}
 \operatorname{tr}r(g)&=t(g),&
 t(g^*)&=t(g),\\
 \bar t(g)&=\sum_i\Psi(g)^{w_i},&
 \bar\chi_{vi}(s)&=\Psi(\operatorname{loc}_vs)^{w_i}.
 \end{aligned}                                                \tag{SF8}
\]
Every trace pairing with the ordinary product is zero
by the defining ordinary determinant relations.
Full matrix span makes the ordinary product itself
zero. No generic automorphy assertion is used in this
descent.

\subsection{The symmetric form and marked normalization}
Equality of the traces of \(r(g)\) and \(r(g^*)^t\),
with full matrix span, gives an intertwiner by
Skolem--Noether. Hence there is an invertible matrix
\(P\) with
\[
             r(g^*)^tP=Pr(g),\qquad
             r(g)^tP\,r(cgc^{-1})=P .
\]
Schur's lemma gives \(P^t=P\) or \(P^t=-P\).
In odd rank the alternating case is impossible
by its determinant and \(2\ne0\).

In rank two the actual symmetric seed excludes it.
For an invertible alternating two-by-two form,
the first displayed relation gives
\[
                    r(g)+r(g^*)=\operatorname{tr}r(g)\,I_2 .
\]
Multiplying by \(r(h)\) and taking traces gives
an identity entirely in the integral function \(t\):
\[
       t(hg)+t(hg^*)=t(g)t(h).                              \tag{SF9}
\]
If the generic form were alternating, (SF9) would
hold in \(C\), and hence at the actual seed.
At its normalized symmetric form \(J\), (SF9) and
full matrix span imply
\(\rho(g)+J\rho(g)^tJ=\operatorname{tr}\rho(g)I_2\).
Its off-diagonal equations are
\(2\rho(g)_{01}=2\rho(g)_{10}=0\).
This contradicts irreducibility of the seed.
Therefore \(P\) is symmetric in both ranks.

Put
\[
 D=r(g_0)-r(g_0^*),\qquad
 A=t(g_0^2)-t(g_0g_0^*)\in C .
\]
Then
\[
       D^tP=-PD,\qquad \operatorname{tr}(D^2)=2A,\qquad
       \bar A=(\Psi(g_0)-\Psi(g_0)^{-1})^2\ne0 .
                                                               \tag{SF10}
\]
Hensel's lemma gives the unique unit \(\zeta\in C\)
with the specified residue and \(\zeta^2=A\).
In ranks two and three, skew-adjointness and (SF10)
give respectively the characteristic polynomials
\[
                       X^2-\zeta^2,\qquad X(X^2-\zeta^2).
\]
For rank three use additionally
\(\det D=-\det D\), hence \(\det D=0\).
Only division by \(2\) is needed.
The distinct eigenvalues are \(\zeta w_i\).
Choose eigenvectors in that order.
Skew-adjointness makes their \(P\)-pairing zero
unless their weights sum to zero.
Nondegeneracy pairs opposite lines, and in rank
three the middle line has nonzero self-pairing.
Scaling these eigenvectors over the algebraically
closed \(K\) puts \(P=J\).
In this basis
\[
 \rho(g^*)=J\rho(g)^tJ,\qquad
 \rho(g_0)-\rho(g_0^*)=\zeta\operatorname{diag}(w_i).
                                                               \tag{SF11}
\]
Conjugation preserves every ordered ordinary product.
The characters of \(\rho\) are still the images of
\(\chi_{vi}\).
Thus \(\rho\) is the actual normalized representation,
with integral trace, integral ordinary characters,
and the integral marked parameter \(\zeta\).

\subsection{The finite entry cone and the full-dimensional point}
The marked projectors in (DC2) are \(C\)-linear
combinations of group matrices.
Consequently
\[
 \rho(g)_{ii}\in C,\quad
 \overline{\rho(g)_{ii}}=\Psi(g)^{w_i},\quad
 \rho(g)_{ij}\rho(h)_{ji}\in\mathfrak m_C\quad(i\ne j).
                                                               \tag{SF12}
\]
For the last assertion insert the two distinct marked
projectors into a trace. Its residual value is zero,
because the residual determinant is the sum of the
distinct diagonal characters.
The same argument with a cyclic list of projectors
puts every closed index cycle in \(C\), and in
\(\mathfrak m_C\) when it has an off-diagonal factor.
The primitive balanced-weight lists of (DC3), with
the polarization reflections, give the same assertion
for every weight-zero entry monomial.

For completeness, all entry spans are finite and
continuous over \(C\). Choose the \(n^2\) group
matrices in (SF6). Solving their trace Gram system
expresses every \(\rho(g)\) with coefficients in
\(\Delta^{-1}C\). Thus their \(C\)-span is a submodule
of a fixed finite free \(C\)-module and is finite.
Every individual entry span is finite as well.
The numerators of these coefficients are continuous
functions of \(g\). Artin--Rees for the finite
matrix span in this free module converts this
continuity into its own \(\mathfrak m_C\)-adic
continuity. The same holds on every entry span.

Let \(D_{\rm gr}\subset K[t,t^{-1}]\) be the graded
\(C\)-algebra generated by
\[
                      t^{w_i-w_j}\rho(g)_{ij}
\]
and the degree-zero ordinary characters.
Its degree-zero part is \(C\).
By (SF12), reduction of this degree-zero part
modulo \(\mathfrak m_C\), and annihilation of every
nonzero-degree part, defines its vertex maximal ideal.
The finite entry spans make \(D_{\rm gr}\) a finite-type
Noetherian \(C\)-algebra.
The fixed-degree topology and the vertex topology
are cofinal by the primitive balanced-monomial
argument of (LS3).
In its vertex completion the displayed matrices
are continuous and have diagonal residual representation.
They therefore receive the universal diagonal
deformation map.
The weight-zero image is contained in the closed
copy of \(C\), and the polynomial entry image lies
in \(D_{\rm gr}\).
Evaluating \(t=1\) gives
\[
             \xi:A_n\longrightarrow K,\qquad
             \beta:B_n\longrightarrow C,\qquad
             \xi|_{B_n}=\iota_C\beta,                       \tag{SF13}
\]
with exactly the matrices \(\rho\) and characters
just constructed.
The copy of \(C\) is closed here also by compactness:
\(C\) is complete Noetherian with finite residue
field, and the vertex completion is Hausdorff.

The two maps \(\beta\Phi_n\) and the quotient
\(\mathcal B_n\to C\) give the same determinant
and the same local characters after injection into
\(K\). They give the same pair over \(C\), by
injectivity, so universality gives
\[
                   \beta\Phi_n=(\mathcal B_n\to C).         \tag{SF14}
\]
It follows that \(\beta\) is surjective.
With \(\mathfrak p=\ker\xi\) and \(Z_B=\ker\beta\),
we have
\[
 \begin{gathered}
 \mathfrak p\text{ Type I},\quad
 M_n\subseteq Z_B=\mathfrak p\cap B_n,\quad
 \dim B_n/Z_B=1+nd .                                       \tag{SF15}
 \end{gathered}
\]
Irreducibility implies Type I by the marked-projector
coordinate-subspace argument.
Equations (SF13)--(SF15) prove the source
\srcid{SeedIn}, including its dimension.
This is the full-dimensional \srcid{SeedOf}
producer; a single arithmetic specialization by
itself would not supply (SF15).

\subsection{The exact support transfers}
For a completed chart \(A_n\to R_x\), pull a prime
\(P\) back to \(\mathcal B_n\) through
\(A_n\leftarrow B_n\leftarrow\mathcal B_n\).
A Type-I prime gives a stable determinant there:
embed \(R_x/P\) into an algebraically closed field.
The marked coordinate projectors and the rank-two
or rank-three zero-entry classification make the
entry representation irreducible, and scalar twisting
and the surjection of arithmetic groups preserve
that property. The determinant at this point is
therefore stable.

For a thick exceptional prime in rank three use
the actual \(s_D\) of (SH7).
Its even entries equal those of \(\rho\), its
middle row off the diagonal is zero, and it has
the two ordered identities (SH9).
Its characteristic polynomials are those of the
outer block plus the middle character, so are
the same as those of \(\rho\).
Twist \(s_D\) by \(\kappa_3\) and compose with
the arithmetic group map. It kills the actual
closed group-algebra kernel:
the Cayley--Hamilton generators are killed by
the matrix Cayley--Hamilton theorem, the ordinary
and starred ordinary generators by the two
identities (SH9), and the inertia generators by
the prime-to-\(\ell\) torsion calculation following
(SF7).

This assertion includes the closure required in
the source. For any target precision, continuity
gives an open normal group subgroup acting trivially
at that precision, and locality of the coefficient
map sends a sufficiently high source maximal-ideal
power into that target precision.
A representative from the source's joint
group/adic closure therefore has zero image at
every target precision. Target separatedness
makes its image zero.
Thus the closed Cayley--Hamilton kernel is
contained in the representation kernel in the
needed quotient direction.

The specialized determinant is not stable,
since its semisimplification has the outer
two-dimensional and middle one-dimensional
summands.
Its closed Cayley--Hamilton algebra is not
semisimple either.
Indeed extend the algebraic closure of its
determinant residue field to an algebraically
closed field containing
\(\operatorname{Frac}(R_x/P)\).
If this finite algebra were semisimple, it
would remain a product of matrix algebras after
this extension. Its representation by \(s_D\)
would split the middle-character quotient,
contrary to the nonsplitting assertion proved
after (SH9).
Here the source field is already an algebraic
closure; this scalar-extension argument is not
an assertion about an arbitrary imperfect field.
Finiteness is the actual
Proposition~\ref{prop:sr-finite-CH}.
We have proved precisely the source
\srcid{IsThick} property.

Now apply Proposition~\ref{prop:common-ordinary-actual-input}
to these stable and thick specializations.
For a selected Type-I prime \(p\) with
\(M_n\subseteq p\cap B_n\), \srcid{P1core} gives
\[
 P\subseteq p,\ P\text{ prime}
                 \quad\Longrightarrow\quad
                          M_n\subseteq P\cap B_n .         \tag{SF16}
\]
For a finite chain of selected exceptional primes
\(p=p_0,\ldots,p_h=q\), with a thick prime below
each successive pair, \srcid{P3core} gives
\[
 \begin{split}
 &\exists P\subseteq p,\quad P\text{ Type I},\
                               M_3\subseteq P\cap B_3\\
 &\hspace{8mm}\Longrightarrow\
 M_3\subseteq P'\cap B_3
       \quad\text{for every Type-I }P'\subseteq q .
 \end{split}                                                \tag{SF17}
\]
The selected chart curves give the actual normalized
\(K[[T]]\)-points, \(N1,N3,N4\), and compatible
determinant maps used by these calls.
Pullback preserves the inclusions in (SF16)--(SF17).
These are exactly \srcid{P1'} and \srcid{P23'}.
Neither transfer asks for the residual outer
constituents to be individually polarized.

\subsection{Meeting the connected characteristic core}
We verify the global incidence assertions as well.
Let
\[
 C_n=\{x\in\operatorname{Spec}A_n:
          x\cap B_n=\mathfrak m_{B_n},\
                          \mathfrak t_+\not\subseteq x\}.
\]
On this fibre the residue-shape equations (LS4)
make every point upper or lower triangular;
the displayed avoidance selects the nonzero
lower-triangular core.
In rank two it is the nonzero strict linear-data
core. In rank three it is the lower-data core.
The explicit polynomial families of
Lemma~\ref{lem:sr-low-interpolation}, and their
actual point identifications in (LS3)--(LS4),
make \(C_n\) connected.
In rank three the nonexceptional part is dense,
using the two actual nonzero strict classes in
Lemma~\ref{lem:frame-strict-classes}.
The fibre is finite type over \(k\), hence
Jacobson. Every nonempty locally closed subset
of it has a closed point. This supplies the
closed-point assertions in
Lemma~\ref{lem:sr-global-propagation}.

Every relevant component meets \(C_n\).
Here is the graded Nakayama argument.
A minimal prime \(\eta\) of \(A_n\) is homogeneous:
its homogeneous core is prime for the ordered
\(\mathbb Z\)-grading and is contained in \(\eta\);
minimality makes them equal.
Suppose
\(\mathfrak t_+\subseteq
             \sqrt{\eta+\mathfrak m_{B_n}A_n}\).
Finite generation gives
\(\mathfrak t_+^m\subseteq\eta+\mathfrak m_{B_n}A_n\)
for some \(m\ge1\).
Every monomial of weight at most \(-2m\)
contains at least \(m\) lower entries, since
each entry weight is at least \(-2\) and
every other entry weight is nonnegative.
Thus for each such weight \(j\), projecting
to the homogeneous component gives
\[
 (A_n)_j\subseteq
          \eta\cap(A_n)_j+\mathfrak m_{B_n}(A_n)_j .
\]
The component is finite over \(B_n\), so
Nakayama makes \((A_n)_j\subseteq\eta\).
For each lower entry its \(2m\)-th power has
such a weight. Primality now puts every lower
entry in \(\eta\), a contradiction on a
relevant component.
There is consequently a prime over
\(\eta+\mathfrak m_{B_n}A_n\) still avoiding
a lower entry.
The Jacobson property on that fibre gives
a closed such point. It lies in \(C_n\).

The full-dimensional seed (SF15) gives a
supported component through this core with
enough local dimension.
Indeed there is a homogeneous prime
\(\eta\subset A_n\) with
\[
 \eta\cap B_n=Z_B,\qquad
 \eta\text{ Type I},\qquad
       \operatorname{trdeg}_{B_n/Z_B}(A_n/\eta)=1 .          \tag{SF18}
\]
To construct it, localize \(A_n/Z_BA_n\)
at \(B_n/Z_B\setminus\{0\}\), take a prime,
and then its homogeneous core.
Its contraction is zero on this field,
so pulling back gives contraction \(Z_B\).
Type I follows from the seed \(\mathfrak p\):
a positive and a negative surviving homogeneous
entry have a balanced product in \(B_n\)
not in \(\mathfrak p\cap B_n=Z_B\).
Neither entry can vanish at \(\eta\).
Pair a surviving odd entry with an entry of
the opposite weight sign to get the same
conclusion for the odd ideal in rank three.
A surviving homogeneous element of nonzero
weight is transcendental over \(B_n/Z_B\),
since a polynomial in it has distinct weights.
Conversely, the balanced products of opposite
weights express every other entry algebraically
in terms of one such entry over
\(\operatorname{Frac}(B_n/Z_B)\).
This proves the transcendence-degree equality.

Choose a closed \(x\in C_n\) over \(\eta\),
as just proved. The exact chart dimension
formula (LS6) gives a prime \(P\) over \(\eta\)
with
\[
        \dim R_x/P=\dim B_n/Z_B+1=nd+2 .
\]
It is supported and Type I.
This is at least the local propagation
threshold \(nd+1\), so it supplies the
supported component required by
Lemma~\ref{lem:sr-global-propagation}.
It is not an inference from a zero-dimensional
arithmetic point.

\subsection{Rank two first, then the exceptional ideal, then rank three}
We now apply the proved local and global
propagation lemmas with these actual objects.

In rank two take the sharp ideal to be zero.
The presentation (LS10), the connectedness
statement (SH4), and the triangular bound (SH1)
make its minimal-component graph connected
at threshold \(2d+1\), and every minimal
component in each core chart is Type I.
The selected primes are (SH3), the supported
locus is \(V(M_2A_2)\), and its support
implication is (SF16).
Lemma~\ref{lem:sr-local-propagation} proves
the chart propagation.
The flat chart (LS5), the core incidence
just checked, its connectedness, and the
full-dimensional seed then give, by
Lemma~\ref{lem:sr-global-propagation},
\[
  M_2\subseteq P_2\cap B_2
                \quad\text{for every rank-two Type-I prime }P_2.
                                                               \tag{SF19}
\]

Only now construct the rank-three exceptional
support ideal.
Use the actual outer-corner map (DC6) and
define \(I_3\) exactly as in (DC8).
Equations (SF4) and (SF19) are the two
rank-two premises (DC7).
Thus the proved corner calculation gives
\[
 \mathfrak m_{B_3}\subseteq\sqrt{I_3+\mathfrak l_3},
 \qquad
 I_3\subseteq P\cap B_3
                  \quad\text{at every rank-three Type-E prime }P .
                                                               \tag{SF20}
\]
This use of the rank-two ideal is prior to
rank-three propagation; no rank-three
automorphy is assumed in obtaining it.

In rank three use the actual sharp ideal
\(I^\sharp\) of (SH5).
Its minimal-component graph is connected
at threshold \(3d+1\), its components have
dimension at least \(3d+2\), and its
dichotomy is Type I or the explicit thick
representation (SH7).
For the ordinary supported prime-selection
step use (SF4); for the thick exceptional
step use (SF20).
The bounds (SH1)--(SH2) and the actual curve
producer (SH3) then supply every selected
prime required in the local propagation
lemma. Its two support implications are
(SF16) and (SF17).
It follows that every Type-I minimal
component through a supported Type-I
component in a core chart is supported.

For a closed core point outside the
exceptional locus the sharp ideal is zero,
so every local minimal component is Type I.
Such a point lies on a relevant global
component. The flat charts also show that
support passes between any two relevant
global components meeting at a closed
core point. The connected core, its dense
nonexceptional locus, its Jacobson
closed points, and its full-dimensional
supported seed are exactly the hypotheses
of Lemma~\ref{lem:sr-global-propagation}.
We conclude
\[
 M_3\subseteq P\cap B_3
                 \quad\text{for every rank-three Type-I prime }P .
                                                               \tag{SF21}
\]
Moreover \(M_3\cap I_3\) has the finite-weight
property and is contained in \(P\cap B_3\)
for every off-triangular prime:
use (SF21) off the exceptional locus and
(SF20) on it. The finite-weight property
passes to this intersection by testing
primes, each of which contains one factor.
This is the actual \srcid{OffFin} conclusion.

\begin{srproposition}[The instantiated shared rank-two/rank-three endpoint]
\label{prop:actual-common-two-three-endpoint}
For the constructed packets, the fixed Hecke ideals
\(J_2,J_3\) satisfy the source's exact
\srcid{ModIn} conditions and give
\(\srcid{GoodPt}_2\), \(\srcid{GoodPt}_3\), and
\(\srcid{FinPt}_3\).
In particular every normalized lift over a complete
coefficient DVR that is irreducible after injection
into a field gives a determinant-pair map from
\(\mathcal B_n/J_n\) with its actual ordered local
characters.
\end{srproposition}

\begin{proof}
Finiteness and dimension are (SF3)--(SF4);
the distinguished unramified condition is (CA14);
the all-point \(N4\) condition is its scalar
residual-Frobenius matrix calculation.
The actual common module proves
\srcid{P1core}, \srcid{P3core}, and
(SF16)--(SF17). The actual seeds give
(SF5)--(SF15). These are all the displayed
fields of \srcid{ModIn}.

For a DVR lift, the actual point map
\(A_n\to A\) of (LS4), followed by a field
injection, is Type I by the marked-projector
classification. Equations (SF19) or (SF21)
make the composite
\(\mathcal B_n\to B_n\to A\) kill \(J_n\).
The determinant and all ordered local
characters agree after the field injection
and hence over \(A\). This is
\srcid{GoodPt}_n.

Finally \srcid{OffFin} and the graded
finite-part argument of
Lemma~\ref{lem:sr-graded-finite-part-criterion}
give \srcid{FinPt}_3.
Explicitly its conclusion is that any
characteristic-\(\ell\) normalized DVR lift
whose ordinary characters are trivial on
\(I_v\cap\ker\Psi\) has, for every \(g\),
the characteristic polynomial
\[
       \prod_{i=0}^{2}
          \bigl(X-\Psi(g)^{w_i}\bigr)
\]
after its specified residue-field embedding.
The constant-character triangularity and
diagonal-character statements used in that
graded argument have their exact scope in
Lemma~\ref{lem:sr-constant-character-triangular}.
This proves the three conclusions.
\end{proof}

This is the completed common support endpoint.
Applying it to an arithmetic characteristic-zero
point still uses its specified algebraic weights
and the classical realization/descent statements.
The original-level local conclusion
\srcid{Th.Unl} is supplied by the separate
crystalline-component analysis, not by replacing
it with the existence of an ordinary flag.

\paragraph{Exact frozen-source provenance.}
The generic trace and form construction is
\srcid{CurveDVRSpecialization.SeedOf}, especially
\srcid{span\_of\_spec}, \srcid{ker\_le},
\srcid{untwist}, and \srcid{seedOf},
together with
\srcid{FrobeniusSmoothVectors.OfTrace.of\_trace}
and \srcid{OfForm.of\_form}.
The finite entry cone is \srcid{PointField.point\_field}.
The core-incidence arguments are
\srcid{MinHom.min\_hom},
\srcid{MonMemPow.l41}, and \srcid{FibreClosed.jac};
the homogeneous seed lift and chart dimension are
\srcid{FrobeniusCornerLiftTwoToThree} and \srcid{Seedpt.seedpt}.
The support transfers are the actual
\srcid{FirstPrimeConditionTransfer.p1Of},
\srcid{StableOf.stableOf},
\srcid{ThickOf.thickOf}, and
\srcid{SecondThirdPrimeConditionTransfer.%
rankThree\_deformationSetup\_splitting\_condition}.
Finally the order (SF19), (SF20), (SF21) is
\srcid{ComponentsGoodRankTwo},
the outer-corner producer \srcid{l53}, and
\srcid{PropagationTheoremGL3.components\_good\_rank\_three}.
Their assembly is
\srcid{CurveDVRSpecialization.AbsMain.absMain}.
All paths lie under
\srcid{Development.AutomorphicArithmetic.AutomorphicVarieties}
at the frozen tag \srcid{alpoge-source-7cbf7a0c}.
